\begin{table}[htbp]
\centering
\caption{County-level border effect, restricted to Colorado and New Mexico}
\label{tab:county_border_capped_conm}
\footnotesize
\begin{tabular}{lccccc}
\toprule
& \multicolumn{2}{c}{Border} & \multicolumn{2}{c}{Border $\times$ Capped} & \\
Outcome & Coef. (\%) & $p$ (cl.) & Coef. (\%) & $p$ (cl.) & $n$ \\
\midrule
Revenue per store & $+0.083$ ($+8.6\%$) & 0.683 & $+0.383$ ($+46.7\%$) & 0.193 & 643 \\
Dispensary density & $+0.613$ ($+84.6\%$) & 0.027 & $-0.916$ ($-60.0\%$) & 0.054 & 643 \\
\bottomrule
\end{tabular}
\begin{minipage}{0.95\textwidth}
\vspace{0.3em}
\footnotesize
\textit{Notes:} Same specification as Table~\ref{tab:county_border_capped_main}, restricted to Colorado and
New Mexico -- the only two states with a genuinely time-varying (not single-snapshot) dispensary count at
this frequency. Montana and Oregon's dispensary counts are a fixed current-year snapshot repeated across every
period, which adds pure noise (not signal) to a quarterly panel; excluding them substantially tightens
dispensary density's estimate (from $p=0.281$ pooled across all 4 states to $p=0.027$ here). County-clustered
standard errors.
\end{minipage}
\end{table}
